\documentclass[11pt]{amsart}

\usepackage[margin=1.08in]{geometry}
\usepackage{amsmath,amssymb,amsthm}
\usepackage{graphicx}
\usepackage[hidelinks]{hyperref}

\newtheorem{theorem}{Theorem}
\newtheorem{lemma}{Lemma}
\theoremstyle{remark}
\newtheorem{remark}{Remark}

\newcommand{\E}{\mathbb E}
\newcommand{\Pp}{\mathbb P}
\newcommand{\R}{\mathbb R}
\newcommand{\N}{\mathbb N}
\newcommand{\eps}{\varepsilon}
\newcommand{\PsiF}{\Psi}
\newcommand{\PhiF}{\Phi}

\title{Convex concentration forces the Hoffmann--Jorgensen Orlicz condition}
\author{Agent lane 06}
\date{Run \texttt{fa\_banach\_001}, 2026-06-29}

\begin{document}
\maketitle

\section*{Status}

This packet gives a full answer to Remark 4.9 of Adamczak--Kutek,
\emph{On Orlicz spaces satisfying the Hoffmann--Jorgensen inequality}
\cite{AK2023}.  They prove that condition \((\mathrm{HJ})\) implies their
convex-Lipschitz concentration inequality (4.9), and ask whether
\((\mathrm{HJ})\) is also necessary when \(\Psi(x)\le C e^{Cx}\).  It is
necessary.

\section*{Source question}

The source defines \(\Psi(x)=e^{\psi(x)}-1\) and says that \(\Psi\) satisfies
\((\mathrm{HJ})\) if there is \(K>0\) such that, for all large \(s,u\),
\[
        \psi(su)\le K\bigl(s\log(1+s)+s\psi(u)\bigr).
\]
For \(\Phi(t)=\Psi(t^2)\), Proposition 4.7 proves a concentration inequality
for convex 1-Lipschitz functions under \((\mathrm{HJ})\).  Immediately after
it, Remark 4.9 states:

\begin{quote}
We note that unlike in the case of inequalities presented in Section 4.1 we do
not know if under an additional assumption that \(\Psi(x)\le Ce^{Cx}\), the
condition \((\mathrm{HJ})\) is necessary for the inequality (4.9) to hold.
\end{quote}

\begin{center}
\includegraphics[width=0.95\textwidth]{figures/open_problem_crop.png}
\end{center}

\section*{Theorem}

Let \(\Psi\) be an Orlicz function, write \(\Psi=e^\psi-1\), and set
\(\Phi(t)=\Psi(t^2)\).  Assume that \(\Psi(x)\le A e^{Ax}\) for all large
\(x\), for some \(A<\infty\).  Suppose that there are constants \(a,A_0>0\)
such that for every \(N\), every random vector
\(X=(X_1,\ldots,X_N)\) with independent real coordinates, every convex
1-Lipschitz \(f:\R^N\to\R\), and every \(t>0\),
\begin{equation}\label{eq:assumed-convex}
\Pp\bigl(|f(X)-\E f(X)|\ge t\bigr)
\le
A_0\exp\left(-a\,\frac{t^2}{(\E \max_{i\le N}|X_i|)^2}\right)
+
\frac{A_0}{\Phi\left(a\,t/\|\max_{i\le N}|X_i|\|_\Phi\right)+1}.
\end{equation}
Then \(\Psi\) satisfies \((\mathrm{HJ})\).

\section*{Idea of the proof}

Use the same sparse Bernoulli variables as in the source's necessity proof for
the empirical-process inequalities, but feed them to the convex function
\(x\mapsto \|x\|_2\).  If there are at least \(s\) rare hits, the Euclidean
norm is of size \(\sqrt{s}\).  Meanwhile the maximum coordinate has expectation
of order \(1/\Psi(u)\) and \(\Phi\)-norm of order \(u^{-1/2}\).  Therefore the
second term in \eqref{eq:assumed-convex}, evaluated at \(t\simeq\sqrt{s}\), is
of order \(\exp[-\psi(csu)]\).  The rare-hit event has Poisson probability at
least \(\exp[-C(s\log(1+s)+s\psi(u))]\).  Comparing the upper and lower bounds
forces the \((\mathrm{HJ})\) inequality.

\section*{Proof}

We shall use constants \(C,c\in(0,\infty)\) depending only on the constants in
\eqref{eq:assumed-convex} and on \(\Psi\).  Their values may change from line
to line.

Fix \(u\) large enough that \(\Psi(u)>2\), and fix \(s\ge 1\).  For each
large \(N\), let \(p_N=(N\Psi(u))^{-1}\), let
\(Y_{N,1},\ldots,Y_{N,N}\) be independent Bernoulli variables with
\(\Pp(Y_{N,i}=1)=p_N\), and set
\[
        X_{N,i}=Y_{N,i}-p_N,\qquad
        X^{(N)}=(X_{N,1},\ldots,X_{N,N}).
\]
Let \(K_N=\sum_{i=1}^N Y_{N,i}\), and let \(Z\) be a Poisson random variable
with parameter \(\lambda=1/\Psi(u)\).  Then \(K_N\) converges in total
variation to \(Z\).

Take the convex 1-Lipschitz function \(f_N(x)=\|x\|_2\).  If
\(K_N\ge \lceil 16s\rceil\) and \(N\) is so large that \(p_N\le 1/2\), then
\[
        f_N(X^{(N)})\ge (1-p_N)\sqrt{K_N}\ge 2\sqrt{s}.
\]
Also
\[
        \E f_N(X^{(N)})
        \le \bigl(\E \|X^{(N)}\|_2^2\bigr)^{1/2}
        = \bigl(Np_N(1-p_N)\bigr)^{1/2}
        \le \Psi(u)^{-1/2}\le 1.
\]
Thus, for \(s\ge 1\),
\[
        \{K_N\ge \lceil 16s\rceil\}
        \subset
        \{f_N(X^{(N)})-\E f_N(X^{(N)})\ge \sqrt{s}\}.
\]
Letting \(N\to\infty\), Stirling's lower bound gives
\begin{equation}\label{eq:poisson-lower}
\liminf_{N\to\infty}
\Pp\bigl(f_N(X^{(N)})-\E f_N(X^{(N)})\ge \sqrt{s}\bigr)
\ge
\exp\{-C(s\log(1+s)+s\psi(u))\}.
\end{equation}
Indeed, the single term
\(\Pp(Z=\lceil 16s\rceil)=e^{-\lambda}\lambda^{\lceil 16s\rceil}/\lceil
16s\rceil!\) has this size, since
\(\log\Psi(u)\le \psi(u)\).

We now estimate the two parameters in \eqref{eq:assumed-convex}.  Let
\[
        W_N=\max_{i\le N}|X_{N,i}|.
\]
If no Bernoulli success occurs, then \(W_N=p_N\); otherwise \(W_N\le 1\).  Hence
\[
        \E W_N\le \Pp(K_N\ge 1)+p_N\le Np_N+p_N\le \frac{2}{\Psi(u)}
\]
for \(N\) large.

For the \(\Phi\)-norm, put \(B_N=\mathbf 1_{\{K_N\ge 1\}}\).  Since
\(\Pp(B_N=1)\le 1/\Psi(u)\),
\[
        \E \Phi(\sqrt{u}\,B_N)
        =
        \Pp(B_N=1)\Psi(u)
        \le 1,
\]
and so \(\|B_N\|_\Phi\le u^{-1/2}\).  Since \(W_N\le B_N+p_N\) and
\(\|\mathbf 1\|_\Phi\) is a fixed finite constant, taking \(N\) still larger
if necessary yields
\begin{equation}\label{eq:parameters}
        \E W_N\le \frac{2}{\Psi(u)},\qquad
        \|W_N\|_\Phi\le \frac{C}{\sqrt{u}} .
\end{equation}

Applying \eqref{eq:assumed-convex} to \(f_N\) with \(t=\sqrt{s}\), using
\eqref{eq:parameters}, and letting \(N\to\infty\), we obtain
\begin{multline}\label{eq:upper}
\limsup_{N\to\infty}
\Pp\bigl(f_N(X^{(N)})-\E f_N(X^{(N)})\ge \sqrt{s}\bigr)\\
\le
C\exp\{-c\,s\Psi(u)^2\}
+
C\exp\{-\psi(csu)\}.
\end{multline}
Here the second term uses \(\Phi(r)=\Psi(r^2)\):
\[
        \Phi\left(a\sqrt{s}/\|W_N\|_\Phi\right)
        \ge
        \Psi(csu).
\]

Combining \eqref{eq:poisson-lower} and \eqref{eq:upper}, and writing
\[
        R(s,u)=C\bigl(s\log(1+s)+s\psi(u)\bigr),
\]
gives
\begin{equation}\label{eq:comparison}
        e^{-R(s,u)}
        \le
        C e^{-c s\Psi(u)^2}
        +
        C e^{-\psi(csu)}
\end{equation}
for all large \(s,u\).

We extract \((\mathrm{HJ})\) from this comparison.  If
\(c s\Psi(u)^2\ge R(s,u)+\log(2C)\), then the first term in
\eqref{eq:comparison} is at most \(e^{-R(s,u)}/2\), and hence
\[
        \psi(csu)\le R(s,u)+\log(2C)
        \le C\bigl(s\log(1+s)+s\psi(u)\bigr).
\]

It remains to consider the complementary case
\[
        c s\Psi(u)^2< R(s,u)+\log(2C).
\]
After increasing the constants and taking \(s,u\) large, this implies
\[
        \Psi(u)^2\le C\bigl(\log(1+s)+\psi(u)\bigr).
\]
Since an Orlicz function is convex, increasing, and nonzero, \(\Psi(u)\ge c u\)
for all large \(u\).  Therefore
\[
        u\le C\bigl(\log(1+s)+\psi(u)\bigr).
\]
Using the assumed upper exponential growth of \(\Psi\), equivalently
\(\psi(x)\le Cx+C\) for all large \(x\), we get
\[
        \psi(csu)\le Csu+C
        \le C\bigl(s\log(1+s)+s\psi(u)\bigr).
\]
Thus in every case
\[
        \psi(csu)\le C\bigl(s\log(1+s)+s\psi(u)\bigr)
\]
for all sufficiently large \(s,u\).  Replacing \(s\) by \(s/c\) and adjusting
the constant gives the standard \((\mathrm{HJ})\) condition:
\[
        \psi(su)\le K\bigl(s\log(1+s)+s\psi(u)\bigr)
\]
for all large \(s,u\).  This proves the theorem.

\section*{Verification and Novelty Notes}

The proof uses only the source paper's definitions, the sparse Bernoulli/Poisson
test-array idea already used in its proof of Proposition 4.5, and elementary
estimates on Orlicz norms.  The new observation is to test the convex
concentration inequality with the Euclidean norm.

Local indexes for \texttt{fa\_banach\_001} were searched for arXiv:2310.04163,
Hoffmann--Jorgensen, Orlicz, and convex concentration; no duplicate packet or
attempt was found.  A bounded web search on 2026-06-29 for the exact arXiv id
and for the phrases ``convex-unbounded'', ``Hoffmann-Jorgensen Orlicz
condition'', and ``necessity convex functions'' found only the source paper and
no later exact resolution.

\section*{Human Review Recommendation}

Review the constants in the final case split, especially the passage from the
comparison inequality \eqref{eq:comparison} to the use of the exponential
growth assumption.  The main probabilistic test and the Orlicz scaling appear
robust.

\begin{thebibliography}{9}
\bibitem{AK2023}
R. Adamczak and D. Kutek,
\emph{On Orlicz spaces satisfying the Hoffmann--Jorgensen inequality},
arXiv:2310.04163.
\end{thebibliography}

\end{document}
